\documentclass[12pt,a4paper]{article}
\usepackage[utf8]{inputenc}
\usepackage[spanish]{babel}
\usepackage[T1]{fontenc}
\usepackage{geometry}
\usepackage{listings}
\usepackage{xcolor}
\usepackage{hyperref}
\usepackage{graphicx}
\usepackage{fancyhdr}
\usepackage{booktabs}
\usepackage{array}
\usepackage{longtable}
\usepackage{enumitem}
\usepackage{tabularx}
\usepackage{titlesec}
\usepackage[most]{tcolorbox}
\usepackage{tikz}
\usetikzlibrary{positioning}

\geometry{margin=2.3cm}
\setlength{\headheight}{15pt}
\AtBeginDocument{\shorthandoff{<>}}

% Paleta profesional inspirada en Terraform
\definecolor{tfPurple}{HTML}{7B42BC}
\definecolor{tfPurpleDark}{HTML}{5C2D91}
\definecolor{tfPurpleLight}{HTML}{BFA4EA}
\definecolor{tfLogoDark}{HTML}{4040B2}
\definecolor{tfLogoLight}{HTML}{5C4EE5}
\definecolor{tfInk}{HTML}{1F2937}
\definecolor{tfSlate}{HTML}{4B5563}
\definecolor{tfBg}{HTML}{F5F3FF}

\titleformat{\section}{\Large\bfseries\color{tfPurpleDark}}{\thesection}{0.7em}{}
\titleformat{\subsection}{\large\bfseries\color{tfPurple}}{\thesubsection}{0.6em}{}
\titleformat{\subsubsection}{\normalsize\bfseries\color{tfInk}}{\thesubsubsection}{0.5em}{}
\setlist[itemize]{itemsep=2pt,topsep=4pt}
\setlist[enumerate]{itemsep=2pt,topsep=4pt}

\hypersetup{
		colorlinks=true,
		linkcolor=tfPurpleDark,
		urlcolor=tfPurple,
		pdftitle={Apuntes Terraform de Cero a Master},
		pdfauthor={Alex}
}

% Configuración de listings para código
\lstset{
		basicstyle=\ttfamily\small,
		inputencoding=utf8,
		backgroundcolor=\color{tfBg},
		frame=single,
		rulecolor=\color{tfPurpleLight},
		breaklines=true,
		commentstyle=\color{tfSlate},
		keywordstyle=\color{tfPurpleDark},
		stringstyle=\color{tfPurple},
		showstringspaces=false,
		captionpos=b,
		tabsize=2,
		literate=
		{á}{{\'a}}1 {é}{{\'e}}1 {í}{{\'i}}1 {ó}{{\'o}}1 {ú}{{\'u}}1
		{Á}{{\'A}}1 {É}{{\'E}}1 {Í}{{\'I}}1 {Ó}{{\'O}}1 {Ú}{{\'U}}1
		{ñ}{{\~n}}1 {Ñ}{{\~N}}1 {ü}{{\"u}}1 {Ü}{{\"U}}1
		{├}{{|--}}3 {└}{{+--}}3 {│}{{|}}1 {─}{{-}}1
}

% Encabezados y pies
\pagestyle{fancy}
\fancyhf{}
\fancyhead[L]{\textcolor{tfPurpleDark}{Apuntes Terraform: De Cero a Master}}
\fancyhead[R]{\thepage}
\fancyfoot[C]{\textcolor{tfSlate}{Infraestructura como Código en producción}}

\newtcolorbox{tfnote}{
	colback=tfBg,
	colframe=tfPurpleLight,
	arc=2mm,
	boxrule=0.6pt,
	left=2mm,
	right=2mm,
	top=1mm,
	bottom=1mm
}

\newtcolorbox{tfchecklist}{
	colback=white,
	colframe=tfPurple,
	arc=1.8mm,
	boxrule=0.5pt,
	left=2mm,
	right=2mm,
	top=1mm,
	bottom=1mm,
	title=Checklist de producción,
	fonttitle=\bfseries,
	coltitle=tfPurpleDark
}

% Logo vectorial fiel al isotipo de Terraform (sin imagen externa)
\newcommand{\terraformbrandlogo}{
\begin{tikzpicture}[x=0.11cm,y=-0.11cm]
	% Basado en las proporciones del SVG público del logo (orientación corregida)
	\fill[tfLogoDark] (36.4,20.2) -- (36.4,39.13) -- (52.8,29.67) -- (52.8,10.72) -- cycle;
	\fill[tfLogoLight] (18.2,10.72) -- (34.6,20.2) -- (34.6,39.13) -- (18.2,29.66) -- cycle;
	\fill[tfLogoLight] (0,0.15) -- (0,19.09) -- (16.4,28.56) -- (16.4,9.62) -- cycle;
	\fill[tfLogoLight] (18.2,50.68) -- (34.6,60.15) -- (34.6,41.21) -- (18.2,31.74) -- cycle;
\end{tikzpicture}
}

\title{{\Huge \textbf{Apuntes Terraform: De Cero a Master}}\\
			 {\Large Guía profesional, práctica y escalable}}
\author{DevOps Notes}
\date{\today}

\begin{document}

\hypersetup{pageanchor=false}

\begin{titlepage}
\centering
\vspace*{1.5cm}
{\terraformbrandlogo\par}
\vspace{0.7cm}
{\Huge\bfseries\color{tfPurpleDark} Apuntes Terraform: De Cero a Master\par}
\vspace{0.4cm}
{\Large\color{tfSlate} Guía profesional, práctica y escalable\par}
\vspace{1.2cm}
\begin{tfnote}
\textbf{Objetivo:} dominar Terraform en escenarios reales: cloud, on-premise, Kubernetes, CI/CD, seguridad, gobernanza y operación enterprise.
\end{tfnote}
\vfill
{\large\textbf{DevOps Notes}}\\[0.2cm]
{\normalsize\today}
\end{titlepage}

\tableofcontents
\newpage
\hypersetup{pageanchor=true}

\section{Cómo usar estos apuntes (ruta recomendada)}

\begin{tfnote}
\textbf{Modo de estudio sugerido:} cada sección debe ir acompañada de práctica real. Regla simple: \textit{leer 30\%, practicar 70\%}.
\end{tfnote}

\subsection{Ruta por niveles}
\begin{tabularx}{\textwidth}{>{\bfseries}p{0.20\textwidth} X}
\toprule
Nivel & Objetivo práctico \\
\midrule
Inicial & Crear infra mínima, entender state, usar \texttt{plan/apply/destroy} sin errores. \\
Intermedio & Diseñar módulos reutilizables, backend remoto, CI con \texttt{plan} en PR. \\
Avanzado & Operar multi-entorno/multi-cuenta, seguridad IaC y políticas obligatorias. \\
Master & Diseñar plataforma completa (cloud + on-prem + Kubernetes) con gobernanza y observabilidad. \\
\bottomrule
\end{tabularx}

\subsection{Competencias que debes dominar}
\begin{itemize}
	\item Modelado declarativo y diseño de módulos mantenibles.
	\item Operación segura del state y migraciones sin downtime.
	\item Integración con Kubernetes y GitOps en entornos reales.
	\item Automatización enterprise (CI/CD, policy as code, FinOps, DR).
\end{itemize}

\section{Qué es Terraform y por qué importa}

Terraform es una herramienta de \textbf{Infraestructura como Código (IaC)} creada por HashiCorp.
Permite describir infraestructura en archivos declarativos y aplicarla de forma reproducible.

\subsection{Problemas que resuelve}
\begin{itemize}
		\item Evita cambios manuales no trazables (\textit{clickops}).
		\item Estandariza despliegues entre entornos (dev, qa, prod).
		\item Facilita revisión por \textit{pull request} y auditoría.
		\item Reduce errores humanos y acelera entregas.
\end{itemize}

\subsection{Conceptos base}
\begin{itemize}
		\item \textbf{Provider}: plugin que traduce Terraform a API cloud (AWS, Azure, GCP, Kubernetes, etc.).
		\item \textbf{Resource}: bloque que crea/modifica/elimina un objeto real.
		\item \textbf{Data source}: lee objetos existentes sin crearlos.
		\item \textbf{State}: archivo que guarda el mapeo entre código y recursos reales.
		\item \textbf{Plan}: diferencia entre estado actual y estado deseado.
\end{itemize}

\section{Terraform en on-premise y entornos híbridos}

Sí, Terraform se puede usar \textbf{perfectamente en on-premise}. No está limitado a nube pública.
Terraform funciona contra cualquier plataforma con provider compatible o API integrable.

\subsection{Casos on-premise frecuentes}
\begin{itemize}
		\item \textbf{VMware vSphere}: VMs, datastores, redes virtuales.
		\item \textbf{OpenStack}: redes, compute, volúmenes.
		\item \textbf{Proxmox}: virtualización y recursos de clúster.
		\item \textbf{Kubernetes on-prem}: namespaces, RBAC, Helm, add-ons.
		\item \textbf{DNS/IPAM/firewall}: Infoblox, F5, Palo Alto, Fortinet, etc.
\end{itemize}

\begin{lstlisting}[caption=Ejemplo básico on-prem con vSphere]
terraform {
	required_providers {
		vsphere = {
			source  = "hashicorp/vsphere"
			version = "~> 2.5"
		}
	}
}

provider "vsphere" {
	user                 = var.vsphere_user
	password             = var.vsphere_password
	vsphere_server       = var.vsphere_server
	allow_unverified_ssl = true
}

resource "vsphere_virtual_machine" "app01" {
	name             = "app01"
	resource_pool_id = data.vsphere_resource_pool.pool.id
	datastore_id      = data.vsphere_datastore.datastore.id
	num_cpus         = 2
	memory           = 4096
	guest_id         = "ubuntu64Guest"
}
\end{lstlisting}

\subsection{On-prem + cloud (híbrido) bien diseñado}
\begin{enumerate}
		\item Estados separados por dominio (on-prem network, cloud network, k8s, apps-bootstrap).
		\item Conectividad explícita (VPN/Direct Connect/ExpressRoute + routing + DNS).
		\item Mismos estándares de tags, naming, políticas y revisiones.
		\item Pipeline única con aprobaciones por entorno.
\end{enumerate}

\subsection{Blueprint visual de arquitectura híbrida}
\begin{center}
\begin{tikzpicture}[font=\small, node distance=0.8cm]
	\node[draw=tfPurpleDark, rounded corners, minimum width=4.2cm, minimum height=1cm, fill=tfBg] (ci) {CI/CD + Policy Checks};
	\node[draw=tfPurpleDark, rounded corners, minimum width=4.2cm, minimum height=1cm, fill=tfBg, below=of ci] (tf) {Terraform (stacks separados)};

	\node[draw=tfPurple, rounded corners, minimum width=4.2cm, minimum height=1.1cm, fill=white, left=2.2cm of tf] (onprem) {On-Prem (vSphere/OpenStack)};
	\node[draw=tfPurple, rounded corners, minimum width=4.2cm, minimum height=1.1cm, fill=white, right=2.2cm of tf] (cloud) {Cloud (AWS/Azure/GCP)};
	\node[draw=tfPurple, rounded corners, minimum width=4.2cm, minimum height=1.1cm, fill=white, below=of tf] (k8s) {Kubernetes + Helm + GitOps};

	\draw[->, thick, color=tfPurpleDark] (ci) -- (tf);
	\draw[->, thick, color=tfPurpleDark] (tf) -- (onprem);
	\draw[->, thick, color=tfPurpleDark] (tf) -- (cloud);
	\draw[->, thick, color=tfPurpleDark] (tf) -- (k8s);
\end{tikzpicture}
\end{center}

\begin{tfnote}
\textbf{Regla práctica:} Terraform provisiona la plataforma (infra), mientras GitOps y herramientas de configuración gestionan el ciclo de vida dinámico de aplicaciones.
\end{tfnote}

\section{Instalación, estructura mínima y flujo de trabajo}

\subsection{Instalación rápida}
\begin{lstlisting}[caption=Instalación y verificación]
# Linux (ejemplo con apt y repositorio oficial)
sudo apt-get update && sudo apt-get install -y gnupg software-properties-common curl
curl -fsSL https://apt.releases.hashicorp.com/gpg | sudo apt-key add -
sudo apt-add-repository "deb [arch=amd64] https://apt.releases.hashicorp.com $(lsb_release -cs) main"
sudo apt-get update && sudo apt-get install terraform

terraform -version
\end{lstlisting}

\subsection{Ciclo de vida básico (imprescindible)}
\begin{lstlisting}[caption=Comandos esenciales]
terraform init      # Descarga providers y prepara backend
terraform fmt       # Formatea el código
terraform validate  # Valida sintaxis y coherencia
terraform plan      # Muestra cambios previstos
terraform apply     # Aplica los cambios
terraform destroy   # Elimina lo gestionado
\end{lstlisting}

\subsection{Estructura mínima recomendada}
\begin{lstlisting}[caption=Estructura de proyecto]
terraform-project/
├── main.tf
├── variables.tf
├── outputs.tf
├── providers.tf
├── versions.tf
└── terraform.tfvars
\end{lstlisting}

\section{HCL a fondo (HashiCorp Configuration Language)}

\subsection{Bloques y atributos}
\begin{lstlisting}[caption=Ejemplo base en HCL]
terraform {
	required_version = ">= 1.6.0"
	required_providers {
		aws = {
			source  = "hashicorp/aws"
			version = "~> 5.0"
		}
	}
}

provider "aws" {
	region = var.aws_region
}

resource "aws_s3_bucket" "logs" {
	bucket = "mi-bucket-logs-unico-12345"
	tags = {
		Project = "plataforma"
		Owner   = "devops"
	}
}
\end{lstlisting}

\subsection{Variables, tipos y validaciones}
\begin{lstlisting}[caption=Variables tipadas y validadas]
variable "aws_region" {
	description = "Región de AWS"
	type        = string
	default     = "eu-west-1"
}

variable "instance_count" {
	description = "Número de instancias"
	type        = number
	default     = 2

	validation {
		condition     = var.instance_count >= 1 && var.instance_count <= 10
		error_message = "instance_count debe estar entre 1 y 10."
	}
}

variable "common_tags" {
	type = map(string)
	default = {
		Environment = "dev"
		Team        = "platform"
	}
}
\end{lstlisting}

\subsection{Locals y funciones}
\begin{lstlisting}[caption=Locals para evitar duplicación]
locals {
	name_prefix = "${var.project}-${var.environment}"
	merged_tags = merge(var.common_tags, {
		ManagedBy = "terraform"
	})
}

# Ejemplos de funciones útiles
# lower(), upper(), format(), join(), split(), lookup(), try(), can(), coalesce()
\end{lstlisting}

\section{State: el corazón de Terraform}

\subsection{Qué contiene el state}
El \texttt{terraform.tfstate} guarda IDs reales, dependencias y metadatos. Si se pierde o corrompe, Terraform puede intentar recrear recursos que ya existen.

\subsection{Buenas prácticas críticas}
\begin{itemize}
		\item Nunca comitear state local en Git.
		\item Usar backend remoto con bloqueo (S3 + DynamoDB, Azure Storage, GCS, Terraform Cloud).
		\item Hacer copias de seguridad y control de acceso estricto.
\end{itemize}

\subsection{Backend remoto en AWS (ejemplo de producción)}
\begin{lstlisting}[caption=Backend S3 con locking en DynamoDB]
terraform {
	backend "s3" {
		bucket         = "tfstate-company-prod"
		key            = "network/prod/terraform.tfstate"
		region         = "eu-west-1"
		dynamodb_table = "terraform-locks"
		encrypt        = true
	}
}
\end{lstlisting}

\section{Dependencias, meta-argumentos y control avanzado}

\subsection{count vs for\_each}
\begin{itemize}
		\item \textbf{count}: útil para colecciones indexadas simples.
		\item \textbf{for\_each}: preferido para colecciones con clave estable (menos drift al cambiar orden).
\end{itemize}

\begin{lstlisting}[caption=for\_each recomendado]
variable "subnets" {
	type = map(object({
		cidr = string
		az   = string
	}))
}

resource "aws_subnet" "this" {
	for_each          = var.subnets
	vpc_id            = aws_vpc.main.id
	cidr_block        = each.value.cidr
	availability_zone = each.value.az

	tags = {
		Name = "subnet-${each.key}"
	}
}
\end{lstlisting}

\subsection{Meta-argumentos clave}
\begin{longtable}{p{0.22\textwidth} p{0.70\textwidth}}
\toprule
\textbf{Meta-argumento} & \textbf{Uso principal} \\
\midrule
\texttt{depends\_on} & Fuerza dependencia explícita cuando Terraform no la infiere. \\
\texttt{lifecycle} & Controla creación/destrucción y evita borrados accidentales. \\
\texttt{provider} & Permite seleccionar provider/alias específico. \\
\texttt{count} & Repite recurso por número. \\
\texttt{for\_each} & Repite recurso por claves de mapa/set. \\
\bottomrule
\end{longtable}

\begin{lstlisting}[caption=lifecycle para minimizar downtime]
resource "aws_launch_template" "app" {
	name_prefix   = "app-"
	image_id      = data.aws_ami.ubuntu.id
	instance_type = "t3.micro"

	lifecycle {
		create_before_destroy = true
	}
}
\end{lstlisting}

\section{Módulos: de básico a diseño enterprise}

\subsection{Qué es un módulo}
Un módulo es una unidad reutilizable de infraestructura. Todo código Terraform está en un módulo; el directorio raíz es el \textit{root module}.

\subsection{Estructura recomendada de módulo}
\begin{lstlisting}[caption=Estructura de módulo reutilizable]
modules/
└── vpc/
		├── main.tf
		├── variables.tf
		├── outputs.tf
		└── README.md
\end{lstlisting}

\subsection{Consumo de módulo}
\begin{lstlisting}[caption=Uso de módulo local]
module "network" {
	source       = "./modules/vpc"
	project      = var.project
	environment  = var.environment
	vpc_cidr     = "10.20.0.0/16"
	subnet_cidrs = ["10.20.1.0/24", "10.20.2.0/24"]
}

output "vpc_id" {
	value = module.network.vpc_id
}
\end{lstlisting}

\subsection{Principios de módulos de calidad}
\begin{enumerate}[label=\arabic*.]
		\item Interfaz clara (variables y outputs mínimos).
		\item Encapsulación (sin exponer detalles internos innecesarios).
		\item Versionado semántico y changelog.
		\item Tests y ejemplos de uso.
		\item Compatibilidad retroactiva en releases mayores/menores.
\end{enumerate}

\section{Workspaces, entornos y estrategia multi-ambiente}

\subsection{Cuándo usar workspaces}
Útiles para diferencias ligeras entre entornos. Para diferencias estructurales grandes, preferir directorios/separación por stacks.

\begin{lstlisting}[caption=Comandos de workspaces]
terraform workspace list
terraform workspace new dev
terraform workspace new prod
terraform workspace select prod
\end{lstlisting}

\subsection{Patrón recomendado en equipos grandes}
\begin{itemize}
		\item Separar por capas: \texttt{network}, \texttt{security}, \texttt{data}, \texttt{apps}.
		\item Un state por stack y entorno.
		\item Promoción por PR y pipeline, nunca cambios manuales en prod.
\end{itemize}

\section{Seguridad real en Terraform}

\subsection{Gestión de secretos}
\begin{itemize}
		\item No guardar secretos en variables planas ni en \texttt{.tfvars} versionados.
		\item Integrar con Secrets Manager, Vault o SSM Parameter Store.
		\item Marcar outputs sensibles.
\end{itemize}

\begin{lstlisting}[caption=Output sensible]
output "db_password" {
	value     = random_password.db.result
	sensitive = true
}
\end{lstlisting}

\subsection{Políticas y compliance}
\begin{itemize}
		\item \textbf{tflint}: calidad y convenciones.
		\item \textbf{tfsec}/\textbf{checkov}: seguridad estática IaC.
		\item \textbf{OPA/Conftest} o Sentinel: políticas obligatorias (guardrails).
\end{itemize}

\section{Testing y calidad profesional}

\subsection{Niveles de validación}
\begin{enumerate}
		\item \textbf{Formato y sintaxis}: \texttt{terraform fmt -check}, \texttt{terraform validate}.
		\item \textbf{Análisis estático}: \texttt{tflint}, \texttt{tfsec}, \texttt{checkov}.
		\item \textbf{Tests de integración}: Terratest (Go) o pipelines efímeros.
\end{enumerate}

\begin{lstlisting}[caption=Pipeline mínima de calidad]
terraform fmt -check -recursive
terraform init -backend=false
terraform validate
tflint
tfsec .
\end{lstlisting}

\section{Import, refactor y migraciones sin downtime}

\subsection{Importar recursos existentes}
\begin{lstlisting}[caption=Import básico]
# Declaras primero el bloque resource
resource "aws_s3_bucket" "legacy" {
	bucket = "bucket-existente"
}

# Luego importas su estado
terraform import aws_s3_bucket.legacy bucket-existente
\end{lstlisting}

\subsection{Refactor seguro con moved blocks}
\begin{lstlisting}[caption=Renombrado sin recrear recursos]
moved {
	from = aws_instance.web
	to   = aws_instance.frontend
}
\end{lstlisting}

\section{CI/CD con Terraform (modelo recomendado)}

\subsection{Flujo empresarial}
\begin{enumerate}
		\item Pull Request ejecuta \texttt{fmt}, \texttt{validate}, seguridad y \texttt{plan}.
		\item Aprobación obligatoria de al menos un revisor.
		\item \texttt{apply} solo en rama protegida y con credenciales de corta duración.
		\item Registro de artefacto \texttt{plan} y auditoría de cambios.
\end{enumerate}

\begin{lstlisting}[caption=Esqueleto CI (pseudo-pipeline)]
stages:
	- lint
	- validate
	- plan
	- apply

plan:
	script:
		- terraform init
		- terraform plan -out=tfplan

apply:
	when: manual
	script:
		- terraform apply -auto-approve tfplan
\end{lstlisting}

\section{Integración profunda con Kubernetes (EKS, AKS, GKE)}

\subsection{Enfoque recomendado: Terraform para plataforma, GitOps para apps}
Terraform destaca creando \textbf{infraestructura de clúster} (red, nodos, IAM, políticas, ingress base, observabilidad base).
Para despliegue continuo de aplicaciones, suele ser mejor GitOps (Argo CD / Flux) para evitar planes gigantes y drift frecuente.

\subsection{Providers típicos en ecosistema Kubernetes}
\begin{itemize}
		\item \textbf{hashicorp/kubernetes}: gestiona recursos Kubernetes desde Terraform.
		\item \textbf{hashicorp/helm}: instala charts Helm de forma declarativa.
		\item \textbf{hashicorp/aws}, \textbf{hashicorp/azurerm}, \textbf{hashicorp/google}: crean clúster administrado y recursos cloud.
\end{itemize}

\begin{lstlisting}[caption=Ejemplo mínimo: EKS + provider kubernetes]
provider "aws" {
	region = var.aws_region
}

data "aws_eks_cluster" "this" {
	name = module.eks.cluster_name
}

data "aws_eks_cluster_auth" "this" {
	name = module.eks.cluster_name
}

provider "kubernetes" {
	host                   = data.aws_eks_cluster.this.endpoint
	cluster_ca_certificate = base64decode(data.aws_eks_cluster.this.certificate_authority[0].data)
	token                  = data.aws_eks_cluster_auth.this.token
}

resource "kubernetes_namespace" "platform" {
	metadata {
		name = "platform"
	}
}
\end{lstlisting}

\subsection{Instalación de add-ons con Helm provider}
\begin{lstlisting}[caption=Ingress NGINX con Helm]
resource "helm_release" "ingress_nginx" {
	name       = "ingress-nginx"
	repository = "https://kubernetes.github.io/ingress-nginx"
	chart      = "ingress-nginx"
	namespace  = "ingress-nginx"

	create_namespace = true
	atomic           = true
	timeout          = 900

	set {
		name  = "controller.replicaCount"
		value = "2"
	}
}
\end{lstlisting}

\subsection{Patrones avanzados en Kubernetes + Terraform}
\begin{enumerate}
		\item Crear clúster y networking con Terraform.
		\item Instalar add-ons base (ingress, cert-manager, external-dns, metrics-server).
		\item Delegar aplicaciones al flujo GitOps.
		\item Separar states: \texttt{cluster}, \texttt{addons}, \texttt{apps-bootstrap}.
\end{enumerate}

\subsection{Errores típicos en k8s con Terraform}
\begin{itemize}
		\item \textbf{Race conditions}: el API server aún no está listo tras crear clúster.
		\item \textbf{Token expirado}: providers que usan auth temporal en pipelines largos.
		\item \textbf{Diff infinito en helm}: valores no deterministas o charts con defaults variables.
		\item \textbf{Drift alto}: si alguien aplica con kubectl/helm fuera de Terraform.
\end{itemize}

\subsection{Kubernetes on-prem (RKE2, K3s, OpenShift, vanilla)}
\begin{itemize}
	\item Terraform puede aprovisionar VMs/red on-prem y bootstrap del clúster.
	\item Para OpenShift, suele coexistir con instaladores nativos + Terraform para infraestructura periférica.
	\item El patrón ganador sigue siendo: \textbf{Terraform para base}, \textbf{GitOps para workloads}.
\end{itemize}

\begin{lstlisting}[caption=Patrón on-prem k8s: namespaces base + helm add-ons]
provider "kubernetes" {
	config_path = "~/.kube/config"
}

resource "kubernetes_namespace" "observability" {
	metadata { name = "observability" }
}

resource "helm_release" "kube_prometheus_stack" {
	name       = "kube-prometheus-stack"
	repository = "https://prometheus-community.github.io/helm-charts"
	chart      = "kube-prometheus-stack"
	namespace  = kubernetes_namespace.observability.metadata[0].name
}
\end{lstlisting}

\section{Arquitectura de plataforma: landing zone y multi-cuenta}

\subsection{Modelo de referencia en cloud pública}
\begin{itemize}
		\item Cuenta/suscripción de \textbf{seguridad} (logs, auditoría, KMS/HSM).
		\item Cuenta/suscripción de \textbf{red compartida} (transit gateway, firewall, DNS).
		\item Cuentas por \textbf{entorno} (dev/qa/prod) y por dominio de negocio.
		\item Identity centralizada con roles de acceso temporal.
\end{itemize}

\subsection{Estrategia por stacks}
\begin{lstlisting}[caption=Separación por stacks y estados]
live/
├── prod/
│   ├── org-security/
│   ├── shared-network/
│   ├── eks-platform/
│   └── apps-bootstrap/
└── dev/
    ├── org-security/
    ├── shared-network/
    ├── eks-platform/
    └── apps-bootstrap/
\end{lstlisting}

\subsection{Terragrunt (opcional, muy usado en enterprise)}
Terragrunt ayuda a reducir duplicación (backend/provider/inputs repetidos), pero añade una capa operativa adicional. Úsalo si tu organización ya tiene muchos stacks y necesidad de DRY extremo.

\subsection{Blueprint de repositorio enterprise}
\begin{lstlisting}[caption=Estructura recomendada para equipos grandes]
platform-iac/
├── modules/
│   ├── network/
│   ├── kubernetes/
│   ├── observability/
│   └── security-baseline/
├── live/
│   ├── prod/
│   │   ├── shared-network/
│   │   ├── eks-platform/
│   │   └── addons/
│   └── dev/
├── policies/
│   ├── opa/
│   └── sentinel/
├── scripts/
├── docs/
└── .github/workflows/
\end{lstlisting}

\section{Patrones de seguridad avanzados}

\subsection{Identidades efímeras en CI/CD}
Evita claves estáticas. Usa OIDC entre tu plataforma CI y cloud provider.

\begin{lstlisting}[caption=Idea de flujo seguro con OIDC]
# 1) CI obtiene token OIDC firmado
# 2) Cloud IAM valida token y emite credenciales temporales
# 3) Terraform usa esas credenciales para plan/apply
\end{lstlisting}

\subsection{Cifrado, secretos y datos sensibles}
\begin{itemize}
		\item Cifrar backend y snapshots de state.
		\item Restringir acceso al state por principio de mínimo privilegio.
		\item Enmascarar variables sensibles en logs de pipeline.
		\item Rotación periódica de secretos y llaves KMS.
\end{itemize}

\subsection{Policy as Code obligatoria}
\begin{itemize}
		\item Reglas de seguridad: no permitir buckets públicos, puertos críticos abiertos, etc.
		\item Reglas de gobierno: naming, tags, regiones permitidas.
		\item Reglas de coste: tamaños máximos de instancias o límites por entorno.
\end{itemize}

\section{Optimización de costes con Terraform}

\subsection{Prácticas de FinOps aplicadas a IaC}
\begin{enumerate}
		\item Etiquetado obligatorio para \texttt{cost-center}, \texttt{owner}, \texttt{environment}.
		\item Ambientes efímeros con TTL automático.
		\item Apagado programado de recursos no productivos.
		\item Revisión continua de derechos de cómputo (rightsizing).
\end{enumerate}

\subsection{Ejemplo de etiquetas obligatorias por variable}
\begin{lstlisting}[caption=Mapa de tags obligatorio]
variable "mandatory_tags" {
	type = map(string)
	validation {
		condition = alltrue([
			contains(keys(var.mandatory_tags), "owner"),
			contains(keys(var.mandatory_tags), "cost-center"),
			contains(keys(var.mandatory_tags), "environment")
		])
		error_message = "Faltan tags obligatorias: owner, cost-center, environment"
	}
}
\end{lstlisting}

\section{Rendimiento y escalabilidad de Terraform}

\subsection{Cuellos de botella habituales}
\begin{itemize}
		\item States gigantes con miles de recursos en un único stack.
		\item Dependencias en cadena que serializan operaciones.
		\item Provider con límites de API y throttling.
\end{itemize}

\subsection{Cómo escalar sin dolor}
\begin{itemize}
		\item Dividir en stacks con ownership claro.
		\item Ejecutar planes en paralelo por dominio cuando sea seguro.
		\item Cachear plugins/providers en runners CI.
		\item Evitar data sources excesivos que consultan APIs en cada plan.
\end{itemize}

\section{Gestión avanzada del state (operaciones delicadas)}

\subsection{Comandos de state que debes dominar con cuidado}
\begin{lstlisting}[caption=State operations]
terraform state list
terraform state show module.network.aws_vpc.main
terraform state mv old.address new.address
terraform state rm aws_s3_bucket.temp
terraform state pull
\end{lstlisting}

\subsection{Reglas de oro antes de tocar state manualmente}
\begin{enumerate}
		\item Haz backup del state.
		\item Ejecuta en ventana controlada y con revisión por pares.
		\item Documenta el cambio en PR/incidencia.
		\item Lanza \texttt{plan} inmediatamente después para validar coherencia.
\end{enumerate}

\section{Provisioners y alternativas modernas}

\subsection{Por qué evitar provisioners cuando puedas}
\texttt{remote-exec} y \texttt{local-exec} son útiles en casos excepcionales, pero aumentan acoplamiento y fragilidad.

\subsection{Alternativas preferibles}
\begin{itemize}
		\item \textbf{Cloud-init}/user-data para bootstrap de máquinas.
		\item Imágenes inmutables (Packer).
		\item Configuración con Ansible/SSM después del aprovisionamiento.
		\item Operadores/controladores en Kubernetes para configuración runtime.
\end{itemize}

\section{Estrategias de migración a Terraform desde infra existente}

\subsection{Plan de migración por fases}
\begin{enumerate}
		\item Inventario de recursos críticos y dependencias.
		\item Import por dominios pequeños (network primero, luego compute/data).
		\item Normalización de naming/tags sin recreación.
		\item Establecer pipeline obligatoria antes de ampliar alcance.
\end{enumerate}

\subsection{Matriz de riesgo recomendada}
\begin{longtable}{p{0.28\textwidth} p{0.22\textwidth} p{0.40\textwidth}}
\toprule
\textbf{Tipo de recurso} & \textbf{Riesgo} & \textbf{Mitigación} \\
\midrule
Red core (VPC/VNet, rutas) & Alto & Entorno réplica, revisión senior, ventana de cambio. \\
Bases de datos & Alto & Backups, réplica, pruebas restore, plan de rollback. \\
Compute stateless & Medio & Despliegue blue/green y health checks. \\
Buckets/objetos & Medio & Versioning, bloqueo de borrado, políticas de retención. \\
\bottomrule
\end{longtable}

\section{Terraform Cloud/Enterprise y operación colaborativa}

\subsection{Cuándo usar Terraform Cloud}
\begin{itemize}
		\item Equipos grandes que requieren ejecución remota centralizada.
		\item Política de aprobación, colas de ejecución y control RBAC unificado.
		\item Integración nativa con policy checks y run tasks.
\end{itemize}

\subsection{Conceptos clave}
\begin{itemize}
		\item \textbf{Workspace}: unidad de ejecución y state.
		\item \textbf{Run}: pipeline de plan/apply.
		\item \textbf{Variable sets}: variables compartidas por múltiples workspaces.
		\item \textbf{Policy checks}: bloqueo de cambios que violen reglas.
\end{itemize}

\section{Patrones multi-cloud y disaster recovery}

\subsection{Multi-cloud realista}
No todo debe ser multi-cloud activo-activo. Suele funcionar mejor:
\begin{itemize}
		\item Cloud principal para producción.
		\item Cloud secundaria para servicios concretos o contingencia.
		\item Abstracción por módulos para evitar lock-in excesivo sin perder eficiencia.
\end{itemize}

\subsection{DR como código}
\begin{enumerate}
		\item Definir RTO/RPO por servicio.
		\item Versionar infraestructura de failover.
		\item Ejecutar simulacros periódicos (\textit{game days}).
		\item Medir tiempos reales de recuperación y ajustar diseño.
\end{enumerate}

\section{Observabilidad de la infraestructura gestionada por Terraform}

\subsection{Qué observar}
\begin{itemize}
		\item Cambios de recursos por despliegue (auditoría de plan/apply).
		\item Eventos de seguridad y cambios de IAM/políticas.
		\item Coste incremental por release.
		\item SLO de plataforma (latencia API, errores, saturación).
\end{itemize}

\subsection{Stack típico}
\begin{itemize}
		\item Logs/auditoría cloud nativos (CloudTrail, Activity Logs, Audit Logs).
		\item Métricas y alertas (Prometheus/Grafana o servicios gestionados).
		\item Trazabilidad de cambios vía PR + run IDs de Terraform.
\end{itemize}

\section{Anti-patrones que debes evitar}

\begin{enumerate}
		\item Un único state para toda la empresa.
		\item Mezclar infraestructura crítica y apps volátiles en el mismo stack.
		\item Variables sin tipo ni validación.
		\item Uso indiscriminado de \texttt{-target} en producción.
		\item Aplicar en local sin trazabilidad ni aprobación.
		\item Confiar solo en \texttt{terraform destroy} sin políticas de protección.
\end{enumerate}

\section{Laboratorios guiados (de básico a experto)}

\subsection*{Lab 1: Fundamentos}
Crear VPC/VNet + instancia + bucket con tags, outputs y variables tipadas.

\subsection*{Lab 2: Módulos y backend remoto}
Extraer red a módulo, configurar backend remoto con locking, y separar dev/prod.

\subsection*{Lab 3: Kubernetes plataforma}
Levantar EKS/AKS/GKE, instalar ingress + cert-manager con Helm provider.

\subsection*{Lab 4: Seguridad y políticas}
Integrar tflint + tfsec + checkov + OPA en PR, bloquear merges inseguros.

\subsection*{Lab 5: Migración controlada}
Importar infraestructura existente y refactorizar con \texttt{moved} sin downtime.

\section{Checklist de salida a producción (Go-Live)}

\begin{tfchecklist}
\begin{itemize}
	\item Backend remoto cifrado, con locking y backups verificados.
	\item \texttt{fmt}, \texttt{validate}, \texttt{tflint}, \texttt{tfsec/checkov} en PR obligatorios.
	\item Política de aprobaciones activa (mínimo 1--2 revisores según criticidad).
	\item Credenciales efímeras (OIDC), sin secretos estáticos en CI.
	\item Runbook de rollback documentado y probado.
	\item Monitorización/alertas activas para recursos críticos.
	\item Matriz de ownership clara por stack (quién aprueba, quién opera, quién responde).
\end{itemize}
\end{tfchecklist}

\subsection{Runbook mínimo de incidente IaC}
\begin{enumerate}
	\item Congelar applies no urgentes.
	\item Ejecutar \texttt{plan} contra estado actual y capturar evidencia.
	\item Identificar si el problema es drift, permisos, dependencia o provider.
	\item Aplicar mitigación segura (rollback, \texttt{state mv/rm}, import, hotfix en módulo).
	\item Publicar postmortem con acciones preventivas.
\end{enumerate}

\section{Troubleshooting experto}

\subsection{Errores comunes y solución}
\begin{longtable}{p{0.36\textwidth} p{0.56\textwidth}}
\toprule
\textbf{Problema} & \textbf{Acción recomendada} \\
\midrule
Drift entre cloud y código & Ejecutar \texttt{plan}, revisar cambios manuales, reimportar si procede. \\
State bloqueado & Verificar proceso activo y liberar lock de forma segura. \\
Recreación inesperada de recurso & Revisar atributos \textit{ForceNew}, usar \texttt{lifecycle} o rediseñar módulo. \\
Dependencias cíclicas & Romper ciclo con outputs intermedios o separación de stacks. \\
\bottomrule
\end{longtable}

\subsection{Comandos de diagnóstico}
\begin{lstlisting}[caption=Comandos útiles de depuración]
terraform state list
terraform state show aws_instance.web
terraform providers
terraform graph
TF_LOG=INFO terraform plan
\end{lstlisting}

\section{Buenas prácticas ``de senior a master''}

\begin{enumerate}
		\item Diseña módulos pequeños, versionados y bien documentados.
		\item Evita monolitos Terraform: divide por dominios y ownership.
		\item Usa naming conventions estrictas y etiquetas obligatorias.
		\item Separa identidad/seguridad de workloads.
		\item Automatiza revisiones de seguridad y coste antes del apply.
		\item Minimiza privilegios de credenciales CI/CD (principio de menor privilegio).
		\item Implementa estrategias de rollback y recovery de state.
\end{enumerate}

\section{Roadmap de estudio: de 0 a master}

\subsection*{Fase 1 (Fundamentos, 1-2 semanas)}
\begin{itemize}
		\item HCL, variables, outputs, providers.
		\item \texttt{init/fmt/validate/plan/apply/destroy}.
		\item Proyecto simple: VPC + 1 instancia + bucket.
\end{itemize}

\subsection*{Fase 2 (Intermedio, 2-4 semanas)}
\begin{itemize}
		\item Módulos, \texttt{for\_each}, data sources.
		\item Backend remoto y locking.
		\item Pipeline con plan en PR.
\end{itemize}

\subsection*{Fase 3 (Avanzado, 1-2 meses)}
\begin{itemize}
		\item Diseño multi-cuenta / multi-región.
		\item Seguridad IaC, políticas y compliance.
		\item Terratest, importaciones y migraciones complejas.
\end{itemize}

\subsection*{Fase 4 (Master, continuo)}
\begin{itemize}
		\item Arquitectura de plataforma interna reusable.
		\item Golden modules, guardrails y catálogo de patrones.
		\item Observabilidad de cambios y optimización de coste cloud.
\end{itemize}

\section{Cheat Sheet final}

\begin{lstlisting}[caption=Comandos que debes dominar]
# Inicialización y validación
terraform init
terraform fmt -recursive
terraform validate

# Plan y apply
terraform plan -out=tfplan
terraform apply tfplan

# Estado
terraform state list
terraform state show <address>

# Workspaces
terraform workspace list
terraform workspace select dev

# Limpieza
terraform destroy
\end{lstlisting}

\section{Glosario rápido (términos clave)}

\begin{longtable}{p{0.28\textwidth} p{0.64\textwidth}}
\toprule
\textbf{Término} & \textbf{Definición práctica} \\
\midrule
\texttt{Provider} & Plugin que conecta Terraform con APIs externas. \\
\texttt{Resource} & Objeto gestionado por Terraform (VM, red, bucket, namespace, etc.). \\
\texttt{Data Source} & Lectura de recursos existentes sin crearlos. \\
\texttt{State} & Base de verdad operativa de Terraform sobre recursos desplegados. \\
\texttt{Plan} & Vista previa exacta de cambios antes de aplicar. \\
\texttt{Module} & Componente reutilizable de infraestructura. \\
\texttt{Workspace} & Contexto de estado aislado dentro de una misma configuración. \\
\texttt{Drift} & Desviación entre lo declarado y lo que realmente existe. \\
\texttt{GitOps} & Operación declarativa continua desde Git para workloads en Kubernetes. \\
\bottomrule
\end{longtable}

\section{Conclusión}

Terraform no es solo una herramienta, es una disciplina de ingeniería.
El nivel \textit{master} se alcanza cuando combinas:
\begin{itemize}
		\item diseño modular sólido,
		\item operaciones seguras en equipo,
		\item automatización con CI/CD,
		\item y gobernanza de infraestructura a escala.
\end{itemize}

Si conviertes estos apuntes en práctica real con proyectos y revisiones técnicas, tendrás un dominio profesional de Terraform.

\end{document}
